\section{SVI calibration and the arbitrage conditions}\label{app:svi}

\subsection*{The fit}

Fitting all five raw parameters of \eqref{eq:svi} at once is possible and
unreliable. The objective has local minima, and on the S\&P 500 expiry at
$T = 0.36$ a direct least-squares run in this sample converged to
$\rho = +0.39$, a positive skew, which is the wrong sign for an equity index
and left 26 basis points of error at the money.

The fit used here is the quasi-explicit reduction. The reduction itself
predates the arbitrage-free treatment, having been proposed in a 2009
whitepaper that \citet{MartiniMingone2022} record as its origin; it is their
paper that characterises which reduced parameters correspond to an
arbitrage-free slice, which is what makes a constrained refit well posed.
Substituting
$y = (\logm - m)/s$ turns \eqref{eq:svi} into
%
\begin{equation}
  \totvar(y) = a + (b\,\rho\,s)\,y + (b\,s)\sqrt{y^2 + 1},
  \label{eq:svi-reduced}
\end{equation}
%
which is linear in the three coefficients $(a,\; c = b\rho s,\; d = bs)$. The
outer search therefore has to handle only $m$ and $s$, with the inner problem
solved exactly as a weighted linear least squares at each candidate pair. Two
constraints make the solution a valid slice: $d \geq |c|$, which is $|\rho| \leq
1$, and $a + \sqrt{d^2 - c^2} \geq 0$, which keeps total variance non-negative
everywhere. The unconstrained inner solution usually satisfies both, and the
constrained solver runs only when it does not, which is what keeps the
two-dimensional outer search affordable.

\subsection*{The weighting, and why it is not uniform}

Residuals are measured in total variance, but the quantity that matters is
error in implied volatility, and the two objectives are not the same. On the
fitted S\&P 500 slice at $T = 0.94$, total variance runs from 0.0245 at the
money to 0.1546 at the deepest quoted strike, a factor of 6.3. Least squares
weights by the square of the residual, so an unweighted fit in $\totvar$ is
dominated by the wing at roughly 40 to 1 and will fit the tail at the expense
of the strikes that matter most for a note struck near the money.

Since $d\totvar = 2\sigma T\, d\sigma$, scaling each weight by
$1/(2\sigma T)^2$ converts the objective into squared volatility error while
leaving it linear in the coefficients, which is what the reduction in
\eqref{eq:svi-reduced} requires. That factor multiplies a base weight taken
from quote tightness, the reciprocal of the relative bid-ask, capped at 50 so
that one unusually tight quote cannot dominate its slice.

\subsection*{Enforcing the two conditions}

The butterfly condition \eqref{eq:butterfly} is evaluated after each fit on a
grid scoped to that expiry's own quoted strikes, padded by 0.10 in
log-moneyness and capped at five standard deviations of total volatility. A
slice violating it anywhere on that grid is refitted with the condition
imposed, rather than having the offending wing smoothed afterwards.

The calendar condition is enforced pairwise between adjacent fitted expiries on
the grid they share. Where total variance decreases in maturity at any point,
the later slice is refitted subject to
$\totvar_{\text{later}}(\logm) \geq \totvar_{\text{earlier}}(\logm)$. Expiries
closer than 25 days apart are thinned beforehand, because two nearly coincident
expiries divide independent fit error by a very small $\Delta T$ and manufacture
violations that are an artefact of the spacing rather than a property of the
surface.

Grid scope is not a detail. Enforcing these conditions on a single wide grid
spanning $\logm \in [-1.5, 1.0]$ asks a two-week slice about strikes far beyond
anything it quotes, where it is purely extrapolating, and the calendar
condition then makes that extrapolated wing a floor for every later expiry. The
constraint propagates forward from the slice with the least information rather
than the most, and the degradation compounds across the term structure. Scoped
instead to the strikes each pair of adjacent expiries actually quotes, the
surfaces fit without that pressure, and the audit in \S\ref{sec:validation}
returns zero failures of either condition across all 68 slices.
